
\paragraph{\textbf{Proposition 2.6.}}
\SGARefTarget{sga3:xiv:2.6}
\emph{Let \(S\) and \(\mathfrak g\) be as in \SGARefLink{sga3:xiv:2.4}{2.4}, and let \(a\) be a
quasi-regular section of \(\mathfrak g\). There is at most one Cartan
subalgebra of \(\mathfrak g\) containing \(a\). Such a Cartan
subalgebra exists---that is, \(a\) is a regular section---if and only
if \(a\) satisfies the following condition:
\[
  \mathfrak d=\Nil(a,\mathfrak g)
\]
is locally a direct-factor submodule of \(\mathfrak g\), and
\(\operatorname{ad}(a)\) induces an automorphism of
\(\mathfrak g/\mathfrak d\). In that case \(\mathfrak d\) is the
unique Cartan subalgebra of \(\mathfrak g\) containing \(a\).}

Suppose first that \(a\) is contained in a Cartan subalgebra
\(\mathfrak d\) of \(\mathfrak g\). Then
\(\operatorname{ad}(a)_{\mathfrak g/\mathfrak d}\) is bijective on
every fiber. Since \(\mathfrak g/\mathfrak d\) is locally free of
finite type, it follows that this endomorphism is an automorphism of
\(\mathfrak g/\mathfrak d\). On the other hand, by \SGARefLink{sga3:xiv:2.1}{2.1} the
endomorphism \(\operatorname{ad}(a)_{\mathfrak d}\) is locally
nilpotent. Hence
\[
  \mathfrak d=\Nil(a,\mathfrak g),
\]
which proves both the uniqueness of \(\mathfrak d\) and the
necessity of the stated regularity criterion.

For sufficiency, observe that the hypothesis on \(a\) implies that
the formation of \(\Nil(a,\mathfrak g)\) commutes with base
extension, and in particular with passage to fibers. Thus the fibers
\(\mathfrak d(s)\) of
\(\mathfrak d=\Nil(a,\mathfrak g)\) are Cartan subalgebras of the
\(\mathfrak g(s)\). Moreover, \(\mathfrak d\) is a Lie subalgebra of
\(\mathfrak g\) by \SGARefLink{sga3:xiii:proposition:4:1}{XIII~4.1}. It is therefore a Cartan subalgebra.

\paragraph{\textbf{Corollary 2.7.}}
\SGARefTarget{sga3:xiv:2.7}
\emph{Under the hypotheses of \SGARefLink{sga3:xiv:2.4}{2.4}, let \(\mathfrak d\) be a Cartan
subalgebra of \(\mathfrak g\), let \(a\) be a section of
\(\mathfrak d\) that is regular in \(\mathfrak g\), and let \(u\) be
an automorphism of \(\mathfrak g\). Then \(u\) preserves
\(\mathfrak d\) if and only if \(u(a)\) is a section of
\(\mathfrak d\).}

Indeed, by transport of structure \(u(a)\) is then a regular section
of \(\mathfrak g\), contained in both Cartan subalgebras
\(\mathfrak d\) and \(u(\mathfrak d)\); \SGARefLink{sga3:xiv:2.6}{Proposition~2.6} therefore
shows that these two subalgebras are equal.
